\section{Discussion and limitations}
The central result is that recovery gain depends on the operating point, not that one training duration is intrinsically preferred. The symmetric intervention makes this distinction falsifiable. Specialization predicts that moving fine-tuning from $3.84$ to $10.24$\,s should increase the long-minus-short recovery interaction. Instead, both fine-tunes gain more at the longer evaluation duration and the interaction becomes smaller after training at $10.24$\,s.

A useful working hypothesis is therefore that fine-tuning changes the model in ways whose measurable benefit is expressed differently across inference conditions. The domain results support that interpretation at a descriptive level. AudioCaps shows the largest recovery and duration interaction, Clotho retains much of both, and hip-hop shows weak recovery with essentially no duration interaction. The public dense text-fine-tuned reference also favors the longer evaluation duration, although it is not recipe-matched. These observations do not identify a mechanism, but they argue against treating duration dependence as a simple memory of the training duration.

The dense diagnostics sharpen the causal boundary. Raw and EMA baselines differ by at most about $0.02$ CLAP, yet AdamW at $10^{-4}$ for $20{,}000$ steps drives every dense gain below $-0.16$ while improving the pruned checkpoint. We did not test a smaller learning rate, so an overly aggressive update for an already converged dense model remains a plausible explanation. The experiment therefore characterizes this short recipe, not dense recovery in general, and it cannot replace the unavailable dense checkpoint after the released $10^{6}$-step recovery.

The remaining limits concern generality and measurement. The study covers one model family, two pruning severities and a small domain grid. Severity~1 is sensitive to prompt selection, although the severity~2 split-half analysis is stable and Clotho independently reproduces a positive duration interaction. The step from $10.24$ to $15.36$\,s is too uncertain to establish a plateau. The argument that short-duration base behavior does not explain the interaction is limited to CLAP and should not be read as a perceptual claim. Human-CLAP, KL and PANNs provide corroboration, but the small author listening exercise remains descriptive only.

The practical implication is simple. A recovery claim should report the paired increment from P to \PFT{} at the intended operating point and at deliberately displaced points. Common generation noise improves precision, while dense, chance or real-data anchors make the scale interpretable. If a mechanism is proposed, the most informative next experiment is one that changes its directional prediction. Here, matching the fine-tuning budget while changing only training duration was more informative than another aggregate score.

\section{Conclusion}
Recovery fine-tuning gain is strongly operating-point dependent in the AudioLDM-M checkpoints studied here. The duration interaction is robust at severity~2 and appears on an independent Clotho prompt set, while a symmetric intervention shows that it does not follow training duration at a matched $20{,}000$-step budget. Recovery remains strong on Clotho but weak on hip-hop, where the duration interaction also disappears despite informative dense anchors. Reporting recovery as a paired function of operating point therefore gives a more faithful account of compressed-model behavior than a single post-fine-tuning score.

\clearpage
{\setlength{\itemsep}{0pt}\setlength{\parsep}{0pt}
\begin{thebibliography}{99}
\bibitem{fang2023diffpruning} G.~Fang, X.~Ma, and X.~Wang, ``Structural pruning for diffusion
models,'' in \emph{Proc.\ NeurIPS}, 2023.
\bibitem{kim2024bksdm} B.-K.~Kim, H.-K.~Song, T.~Castells \emph{et al.}, ``BK-SDM: A lightweight, fast,
and cheap version of Stable Diffusion,'' in \emph{Proc.\ ECCV}, 2024.
\bibitem{fang2025tinyfusion} G.~Fang, K.~Li, X.~Ma, and X.~Wang, ``TinyFusion: Diffusion transformers
learned shallow,'' in \emph{Proc.\ CVPR}, 2025.
\bibitem{liu2023audioldm} H.~Liu, Z.~Chen, Y.~Yuan \emph{et al.}, ``AudioLDM: Text-to-audio generation with latent diffusion models,'' in
\emph{Proc.\ ICML}, 2023.
\bibitem{singh2026pruning} A.~Singh, Y.~Yuan, Y.~Chen, W.~Wang, and M.~D.~Plumbley, ``Efficient
text-to-audio generation via pruning,'' arXiv:2607.13330, 2026.
\bibitem{wu2023clap} Y.~Wu, K.~Chen, T.~Zhang \emph{et al.},
``Large-scale contrastive language-audio pretraining with feature fusion and keyword-to-caption
augmentation,'' in \emph{Proc.\ ICASSP}, 2023.
\bibitem{huang2023makeanaudio} R.~Huang, J.~Huang, D.~Yang \emph{et al.}, ``Make-An-Audio:
Text-to-audio generation with prompt-enhanced diffusion models,'' in \emph{Proc.\ ICML}, 2023.
\bibitem{ghosal2023tango} D.~Ghosal, N.~Majumder, A.~Mehrish, and S.~Poria, ``Text-to-audio generation
using instruction guided latent diffusion model,'' in \emph{Proc.\ ACM Int.\ Conf.\ Multimedia (MM)},
2023, pp.~3590--3598.
\bibitem{kilgour2019fad} K.~Kilgour, M.~Zuluaga, D.~Roblek \emph{et al.}, ``Fr\'echet Audio
Distance: A reference-free metric for evaluating music enhancement algorithms,'' in
\emph{Proc.\ Interspeech}, 2019.
\bibitem{gui2024fad} A.~Gui, H.~Gamper, S.~Braun \emph{et al.}, ``Adapting Fr\'echet Audio
Distance for generative music evaluation,'' in \emph{Proc.\ ICASSP}, 2024.
\bibitem{chung2025kad} Y.~Chung, P.~Eu, J.~Lee \emph{et al.}, ``KAD: No more FAD!
An effective and efficient evaluation metric for audio generation,'' arXiv:2502.15602, 2025.
\bibitem{kong2020panns} Q.~Kong, Y.~Cao, T.~Iqbal \emph{et al.}, ``PANNs:
Large-scale pretrained audio neural networks for audio pattern recognition,'' \emph{IEEE/ACM
Trans.\ Audio, Speech, Lang.\ Process.}, vol.~28, pp.~2880--2894, 2020.
\bibitem{takano2025humanclap} T.~Takano, Y.~Okamoto, Y.~Kanamori \emph{et al.}, ``Human-CLAP:
Human-perception-based contrastive language-audio pretraining,'' in \emph{Proc.\ APSIPA ASC}, 2025,
pp.~131--136.
\bibitem{li2026finelap} X.~Li, X.~Xu, Z.~Ma \emph{et al.}, ``FineLAP: Taming
heterogeneous supervision for fine-grained language-audio pretraining,'' in \emph{Proc.\ ACL}, 2026,
pp.~10393--10408.
\bibitem{kumar2022finetuning} A.~Kumar, A.~Raghunathan, R.~Jones \emph{et al.}, ``Fine-tuning
can distort pretrained features and underperform out-of-distribution,'' in \emph{Proc.\ ICLR}, 2022.
\bibitem{kim2019audiocaps} C.~D.~Kim, B.~Kim, H.~Lee, and G.~Kim, ``AudioCaps: Generating captions
for audios in the wild,'' in \emph{Proc.\ NAACL-HLT}, 2019.
\bibitem{drossos2020clotho} K.~Drossos, S.~Lipping, and T.~Virtanen, ``Clotho: An audio captioning dataset,'' in \emph{Proc.\ ICASSP}, 2020.
\bibitem{agostinelli2023musiclm} A.~Agostinelli, T.~I.~Denk, Z.~Borsos \emph{et al.}, ``MusicLM:
Generating music from text,'' arXiv:2301.11325, 2023.
\bibitem{song2021ddim} J.~Song, C.~Meng, and S.~Ermon, ``Denoising diffusion implicit models,'' in
\emph{Proc.\ ICLR}, 2021.
\bibitem{wang2026ttabench} H.~Wang, C.~Liu, J.~Chen \emph{et al.}, ``TTA-Bench: A comprehensive benchmark for evaluating text-to-audio models,'' in \emph{Proc.\ AAAI}, 2026, pp.~33512--33520.
\bibitem{bibbo2026repo} G.~Bibb\'o, A.~Singh, and M.~D.~Plumbley, ``Companion repository for operating-point evaluation of recovery fine-tuning,'' 2026. [Online]. Available: \url{https://github.com/gbibbo/audioldm-modality-swap-pruning/tree/main/paper-operating-point-recovery}
\end{thebibliography}}

\end{document}
